\section{Einleitung und Motivation}

\subsection{Problemstellung}
Getrennt erfasste Abfälle können nur dann gezielt weiterverarbeitet werden, wenn Materialarten zuverlässig unterschieden werden. Kamerabasierte Objekterkennung ist dafür ein möglicher Baustein: Sie liefert zu einer Klasse auch die Position eines Gegenstands im Bild. Visuell ähnliche oder verformte Gegenstände, wechselnde Hintergründe und begrenzte Rechenleistung erschweren jedoch den Einsatz. MIRA untersucht dieses Problem zunächst als reproduzierbaren Softwareprototyp; Aussagen über die Sortierleistung einer vollständigen Anlage sind damit noch nicht möglich.

\subsection{Ziel und Forschungsfrage}
Ziel ist eine nachvollziehbare Pipeline für fünf Klassen: \textit{glass, metal, paper, plastic} und \textit{trash}. Im Mittelpunkt steht nicht der Nachweis eines fertigen Sortierroboters, sondern die folgende Forschungsfrage:

\begin{quote}
\emph{Welcher messbare Zielkonflikt zwischen Detektionsgüte, Modellgröße und CPU-Latenz ergibt sich beim Training und bei der INT8-Quantisierung eines kompakten YOLO11n-Modells für fünf Abfallklassen?}
\end{quote}

Die Detektionsgüte wird mit mAP50--95 als Hauptmetrik sowie mAP50, Precision, Recall und dem schwellenwertabhängigen $F_1$-Score beschrieben. Modellgröße und Latenz werden für dasselbe exportierte Modell erfasst. Messungen auf einem Raspberry Pi und die Kopplung an einen Roboterarm bilden eine spätere Projektstufe.

\subsection{Themenfindung}
Die Projektidee entstand aus Beobachtungen der Mülltrennung im schulischen Alltag und bei Besuchen lokaler Recyclinghöfe. Daraus entwickelte sich die technische Frage, wie weit eine kleine, lokal ausführbare Computer-Vision-Pipeline mit dokumentierten Daten und Messgrößen kommt.

\subsection{Abgrenzung der Projektphasen}
Stage~A ist ein abgeschlossener Machbarkeitsnachweis zur Klassifikation einzelner Bilder in vier Klassen. Die späteren EXP-013 bis EXP-017 sind historische Fünfklassen-Detektionsläufe mit unterschiedlichen Datenmischungen. Sie motivieren die aktuelle Untersuchung, ersetzen wegen ihrer unterschiedlichen Datengrundlagen aber keinen kontrollierten Vergleich. Der neu aufgebaute, annähernd ausgeglichene Datensatz und der darauf geplante Lauf bilden den aktuellen Arbeitsstand. Raspberry Pi, Kameraaufbau und Roboterarm sind noch nicht experimentell integriert.

\subsection{Aufbau der Arbeit}
Nach den notwendigen Grundlagen werden Datensätze, Trennregeln und Messgrößen beschrieben. Anschließend folgen Implementierung, bisherige Ergebnisse, deren Grenzen und die geplanten nächsten Versuche.
